\section{Urgency: patients are waiting now}
\label{sec:urgency}

The final appeal is timing. Legislation accelerates when delay is shown to have a
cost: patients are waiting now, lives are being lost now, and the technology is
advancing faster than the rules that govern it. But urgency must be handled with
care. Research on fear appeals finds that they motivate action only when the
audience also believes a solution is within reach; a catastrophe presented without
a credible remedy tends to paralyze rather than move \cite{kok2013finding}. So this
section pairs the pressure of time with a remedy.

The gap is real. Capability, in the form of Physical AI, surgical humanoids, and
autonomous pipelines, is rising quickly, while regulation, built for an earlier
era and supplemented mainly by guidance, lags behind \cite{fda2021samd}. The
distance between the two curves is where preventable harm accumulates.
Figure 11 shows the gap and the single instrument that closes it.

\begin{narrativefig}
\node[nftwo] (c1) at (0,0) {Physical\\AI}; \node[nftwo] (c2) at (4,0) {Surgical\\humanoids};
\node[nftwo] (c3) at (8,0) {Autonomous\\pipelines};
\node[nfone] (r1) at (0,-1.6) {Legacy\\rules}; \node[nfone] (r2) at (4,-1.6) {Guidance\\only};
\node[nfrisk] (r3) at (8,-1.6) {Regulation\\gap};
\node[nfact] (act) at (12, 0) {H. R.\\9510};
\node[nfhope] (g) at (12,-1.6) {Patients\\protected};
\draw[nfedge] (c1)--(c2); \draw[nfedge] (c2)--(c3); \draw[nfedge] (r1)--(r2); \draw[nfedge] (r2)--(r3);
\draw[nfedge] (c3) to[out=0,in=120, looseness=0] (act.west);
\draw[nfedge] (r3) to[out=0,in=-120] (act.west);
\draw[nfedge] (act)--(g);
\end{narrativefig}
\vspace{-0.8cm}
\figcaption{Figure 11. Closing the capability-regulation gap: the Act is the bridge that joins fast-moving capability to
lagging rules.}

\autoref{tab:urgency} sets the cost of another year of delay against the cost of
acting, so the urgency points toward a decision rather than toward despair.

\needspace{9\baselineskip}\begin{table}[H]
\footnotesize
\begin{tabularx}{\textwidth}{@{}L{6.4cm} Y@{}}
\toprule
\textbf{Dimension} & \textbf{Effect of one more year}\\
\midrule
Patients waiting & Trial access delayed for those with the least time\\
Preventable harm & Unverified tools deployed without a gate\\
Public trust & Skepticism hardens without a transparent record\\
Cost of acting & Bounded, authorized, and budget-scored in the bill\\
\bottomrule
\end{tabularx}
\caption{The cost of delay against the bounded cost of acting.}
\label{tab:urgency}
\end{table}

\clearpage

The eight appeals converge on one request. They are not eight separate arguments
but eight doors into the same room, and each one opens onto a specific provision
of the bill, summarized in \autoref{tab:synthesis}. This is the bridge the
advocacy literature describes, where the emotional story makes the data memorable
and the data makes the story actionable \cite{moreland2015hearing, kok2013finding}.

\needspace{11\baselineskip}\begin{table}[H]
\footnotesize
\begin{tabularx}{\textwidth}{@{}L{4.0cm} Y@{}}
\toprule
\textbf{Pillar} & \textbf{Provision of H. R. 9510 that answers it}\\
\midrule
Compassion & A verified pathway a patient can enter, with consent confirmed and recorded\\
Fear of harm & The mandatory ten-gate verification before any action executes\\
Moral outrage & Subgroup bias surveillance, measured and publicly reported\\
Hope & The authorization that funds a tested national trial platform\\
Responsibility & The conforming amendments that make verification the standard of care\\
Protection & Verified informed consent and a per-subgroup outcome audit\\
Trust & Each gate bound to a published standard, with a hash-chained audit trail\\
Urgency & A defined effective date and a budget-scored, bounded cost of acting\\
\bottomrule
\end{tabularx}
\caption{The eight pillars and the specific provisions of H. R. 9510 that answer them.}
\label{tab:synthesis}
\end{table}

We therefore ask the committee to enact H. R. 9510 (v5.0) into Federal law
\cite{Kawchak2026HR9510Bill}. The quiet questions, asked eight times, resolve into
a single answer: we no longer have to carry this task with unaided human attention,
and we should not ask patients to wait while we decide to stop.
